%========================================
% MatinBook — Stage 05: Math System Test
% Version: v1.1
% Purpose:
%   Validate the entire mathematical system of MatinBook v1.1:
%     1. Inline and display equations
%     2. Multi-line equation environments (equation, align, gather)
%     3. Matrices (bracketed, parenthesized, determinant)
%     4. Custom operators (vector calculus, linear algebra, ...)
%     5. Smart delimiters (abs, norm, ceil, floor, inner, paren)
%     6. Set notation (set, setbuilder)
%     7. Derivatives and integrals
%     8. Special functions (Gamma, sinc, erf)
%     9. Number sets (N, Z, Q, R, C)
%
% Compilation:
%   xelatex -shell-escape stage05-math.tex
%   xelatex -shell-escape stage05-math.tex
%   (2 passes required for cover + cross-references)
%
% Expected outcome:
%   A professionally typeset PDF demonstrating that every
%   mathematical feature of MatinBook works correctly, with
%   proper spacing, numbering, and cross-references.
%========================================

%------------------------------------------------------------------------
% Search Path (for tests directory)
%------------------------------------------------------------------------
\makeatletter
\def\input@path{{../}}
\makeatother

%------------------------------------------------------------------------
% Document Class
%------------------------------------------------------------------------
\documentclass{matinbook}

%------------------------------------------------------------------------
% Book Metadata
%------------------------------------------------------------------------
\title{آزمون سیستم ریاضی}
\author{تیم توسعه MatinBook}
\date{مهر ۱۴۰۵}

\booktitle{آزمون سیستم ریاضی}

\renewcommand{\repository}{github.com/matinmahmoudi-cs/matinbook}

%========================================================================
% DOCUMENT
%========================================================================
\begin{document}

%------------------------------------------------------------------------
% FRONT COVER
%------------------------------------------------------------------------
\makecover
    {آزمون سیستم ریاضی}
    {ارزیابی معادلات، ماتریس‌ها، عملگرها و دلیمترها}
    {تیم توسعه MatinBook}
    {مهر ۱۴۰۵}

%------------------------------------------------------------------------
% FRONT MATTER (symmetric margins)
%------------------------------------------------------------------------
\frontmatter
\frontmattergeometry

% Title page
\maketitle

% Copyright / Colophon
\thispagestyle{empty}
\vfill
\begin{center}
    {\large\textbf{شناسنامه آزمون}}

    \vspace{0.8cm}

    \begin{tabular}{rl}
        عنوان: & آزمون سیستم ریاضی \\
        نسخه: & \matinbookversion \\
        تاریخ: & \matinbookdate \\
        موتور: & \lr{XeLaTeX} \\
        مجوز: & \lr{MIT} \\
    \end{tabular}

    \vspace{0.8cm}

    {\small این سند بخشی از مجموعه آزمون‌های یکپارچگی \lr{MatinBook} است.}
\end{center}
\clearpage

% Preface
\chapter*{پیشگفتار}
\addcontentsline{toc}{chapter}{پیشگفتار}

این سند، پنجمین آزمون از مجموعه‌ی آزمون‌های یکپارچگی
\lr{MatinBook} نسخه‌ی \lr{\matinbookversion} است. هدف آن،
اعتبارسنجی کامل سیستم ریاضی کلاس در نه حوزه‌ی متمایز است:
معادلات درون‌خطی و نمایشی، محیط‌های چندخطی، ماتریس‌ها،
عملگرهای سفارشی، دلیمترهای هوشمند، نماد مجموعه‌ها،
مشتق و انتگرال، توابع ویژه و مجموعه‌های عددی.

ریاضیات، ستون فقرات کتاب‌های فنی است. یک سیستم ریاضی
حرفه‌ای باید نمادهای پیچیده را با فاصله‌گذاری دقیق،
شماره‌گذاری منظم و ارجاع هوشمند ارائه دهد.
\lr{MatinBook} با استفاده از \lr{unicode-math}،
\lr{mathtools} و عملگرهای سفارشی، این اهداف را محقق
می‌سازد.

هر بخش از این آزمون، با مثال‌های عملی و ارجاعات متقابل
همراه است تا خواننده بتواند درستی رفتار هر زیرسیستم را
در بافت واقعی یک کتاب ارزیابی کند.

\clearpage

% Table of Contents
\tableofcontents
\clearpage

% List of Figures
\listoffigures
\clearpage

% List of Tables
\listoftables
\clearpage

%------------------------------------------------------------------------
% MAIN MATTER (wide margin for notes)
%------------------------------------------------------------------------
\mainmatter
\mainmattergeometry

%========================================================================
% CHAPTER 1 — BASIC EQUATIONS
%========================================================================
\chapter{معادلات پایه}
\label{chap:equations}

\section{معادلات درون‌خطی}
\label{sec:equations-inline}

ریاضیات درون‌خطی برای اشاره به نمادها و عبارات کوتاه استفاده
می‌شود: $E = mc^{2}$، قضیه‌ی فیثاغورث $a^{2} + b^{2} = c^{2}$،
و اتحاد اویلر $e^{i\pi} + 1 = 0$.

همچنین می‌توان از مجموعه‌های عددی استاندارد استفاده کرد:
$\N$ (طبیعی)، $\Z$ (صحیح)، $\Q$ (گویا)، $\R$ (حقیقی)،
و $\C$ (مختلط). رابطه‌ی شمول میان این مجموعه‌ها در
\cref{chap:number-sets} بررسی خواهد شد.

\mnote{مجموعه‌های عددی با \lr{\textbackslash N}، \lr{\textbackslash Z}، \lr{\textbackslash Q}، \lr{\textbackslash R} و \lr{\textbackslash C} نوشته می‌شوند.}

\section{معادلات شماره‌گذاری‌شده}
\label{sec:equations-numbered}

\begin{equation}
    f(x) = x^{3} - 3x^{2} + 2x - 1
    \label{eq:cubic}
\end{equation}

\begin{equation}
    \frac{d}{dx}\left( \sin x \right) = \cos x
    \label{eq:derivative-sin}
\end{equation}

معادله‌ی \cref{eq:cubic} یک چندجمله‌ای درجه سه است، و
معادله‌ی \cref{eq:derivative-sin} مشتق تابع سینوس را
نشان می‌دهد.

\section{معادلات چندخطی}
\label{sec:equations-multiline}

محیط \lr{align} برای معادلات هم‌تراز:

\begin{align}
    (a + b)^{2} &= (a + b)(a + b) \nonumber \\
    &= a^{2} + ab + ba + b^{2} \nonumber \\
    &= a^{2} + 2ab + b^{2} \label{eq:binomial}
\end{align}

اتحاد دو جمله‌ای در معادله‌ی \cref{eq:binomial} نمایش داده
شده است. محیط \lr{gather} برای معادلات غیرهم‌تراز:

\begin{gather}
    x = a + b + c + d + e + f \label{eq:long1} \\
    y = \alpha + \beta + \gamma + \delta + \epsilon \label{eq:long2}
\end{gather}

معادلات \cref{eq:long1} و \cref{eq:long2} نمونه‌هایی از
محیط \lr{gather} هستند. محیط \lr{align} با چند ستون:

\begin{align}
    x &= a + b & y &= c + d \label{eq:multi1} \\
    u &= e + f & v &= g + h \label{eq:multi2}
\end{align}

%========================================================================
% CHAPTER 2 — MATRICES AND VECTORS
%========================================================================
\chapter{ماتریس‌ها و بردارها}
\label{chap:matrices}

\section{انواع ماتریس}
\label{sec:matrices-types}

ماتریس با کروشه (\lr{\textbackslash mat}):

\[
A = \mat{1 & 2 & 3 \\ 4 & 5 & 6 \\ 7 & 8 & 9}
\]

ماتریس با پرانتز (\lr{\textbackslash pmat}):

\[
B = \pmat{a_{11} & a_{12} \\ a_{21} & a_{22}}
\]

دترمینان (\lr{\textbackslash vmat}):

\[
\det(B) = \vmat{a_{11} & a_{12} \\ a_{21} & a_{22}}
= a_{11}a_{22} - a_{12}a_{21}
\]

ماتریس همانی $3 \times 3$:

\[
I_{3} = \mat{1 & 0 & 0 \\ 0 & 1 & 0 \\ 0 & 0 & 1}
\]

\section{عملگرهای ماتریسی}
\label{sec:matrices-operators}

\begin{itemize}
    \item \textbf{اثر (\lr{Trace}):}
    $\trace(A) = a_{11} + a_{22} + \cdots + a_{nn}$
    \item \textbf{دترمینان:} $\det(A) = \abs{A}$
    \item \textbf{ترانهاده:} $A^{T} = [a_{ji}]$
    \item \textbf{ماتریس قطری:}
    $\diag(d_{1}, d_{2}, \ldots, d_{n})$
    \item \textbf{رتبه:} $\rank(A)$
\end{itemize}

\section{بردارها}
\label{sec:matrices-vectors}

\[
\vec{v} = \mat{v_{1} \\ v_{2} \\ v_{3}}, \quad
\vec{w} = \pmat{w_{1} & w_{2} & w_{3}}
\]

ضرب داخلی: $\inner{\vec{u}}{\vec{v}} = u_{1}v_{1} + u_{2}v_{2} + u_{3}v_{3}$

نرم اقلیدسی: $\norm{\vec{v}} = \sqrt{\inner{\vec{v}}{\vec{v}}}$

%========================================================================
% CHAPTER 3 — SMART DELIMITERS
%========================================================================
\chapter{دلیمترهای هوشمند}
\label{chap:delimiters}

\section{قدر مطلق و نرم}
\label{sec:delimiters-abs-norm}

قدر مطلق ساده: $\abs{-5} = 5$

قدر مطلق با اندازه‌ی دستی: $\abs[\Big]{\frac{a}{b}}$

نرم بردار: $\norm{\vec{x}}$

نرم با اندیس: $\norm{x}_{2} = \sqrt{\sum_{i} x_{i}^{2}}$

\section{سقف و کف}
\label{sec:delimiters-ceil-floor}

\[
\ceil{3.14} = 4, \quad
\floor{3.14} = 3, \quad
\ceil[Big]{\frac{22}{7}} = 4
\]

\section{ضرب داخلی}
\label{sec:delimiters-inner}

\[
\inner{u}{v}, \quad
\inner[\Big]{\frac{x}{2}}{\frac{y}{3}}
\]

\section{مجموعه}
\label{sec:delimiters-set}

\[
A = \set{1, 2, 3, 4, 5}, \quad
B = \setbuilder{x \in \R}{x > 0}
\]

\section{پرانتز هوشمند}
\label{sec:delimiters-paren}

\[
\paren{x + y}, \quad
\paren[\Big]{\frac{a + b}{c + d}}, \quad
\paren*{\frac{\sum_{i=1}^{n} x_{i}}{n}}
\]

%========================================================================
% CHAPTER 4 — OPERATORS AND SPECIAL FUNCTIONS
%========================================================================
\chapter{عملگرها و توابع ویژه}
\label{chap:operators}

\section{عملگرهای حسابان برداری}
\label{sec:operators-vector}

گرادیان تابع اسکالر $f$:

\[
\grad f = \pderiv{f}{x}\hat{i} + \pderiv{f}{y}\hat{j} + \pderiv{f}{z}\hat{k}
\]

کرل میدان برداری $\vec{F}$:

\[
\curl \vec{F} = \nabla \times \vec{F}
\]

دیورژانس میدان برداری $\vec{F}$:

\[
\diver \vec{F} = \nabla \cdot \vec{F}
\]

\section{عملگرهای جبر خطی}
\label{sec:operators-linalg}

\begin{align}
    \rank(A) &= \text{بعد فضای ستونی} \label{eq:rank} \\
    \trace(A) &= \sum_{i=1}^{n} a_{ii} \label{eq:trace} \\
    \diag(d_{1}, \ldots, d_{n}) &=
    \mat{d_{1} & 0 & \cdots & 0 \\
        0 & d_{2} & \cdots & 0 \\
        \vdots & \vdots & \ddots & \vdots \\
        0 & 0 & \cdots & d_{n}} \label{eq:diag}
\end{align}

معادلات \cref{eq:rank}، \cref{eq:trace} و \cref{eq:diag}
عملگرهای اصلی جبر خطی را نشان می‌دهند.

\section{عملگرهای نظریه اعداد}
\label{sec:operators-numbertheory}

\begin{itemize}
    \item بزرگ‌ترین مقسوم‌علیه مشترک: $\gcdop(a, b)$
    \item کوچک‌ترین مضرب مشترک: $\lcm(a, b)$
    \item تابع علامت: $\sgn(x)$
\end{itemize}

\section{بهینه‌سازی}
\label{sec:operators-optimization}

\begin{align}
    \argmax_{x \in \R} f(x) &= \text{مقدار } x \text{ که } f(x) \text{ را بیشینه کند} \\
    \argmin_{x \in \R} f(x) &= \text{مقدار } x \text{ که } f(x) \text{ را کمینه کند}
\end{align}

%========================================================================
% CHAPTER 5 — CALCULUS
%========================================================================
\chapter{حساب دیفرانسیل و انتگرال}
\label{chap:calculus}

\section{مشتق}
\label{sec:calculus-derivatives}

مشتق معمولی:

\[
\deriv{}{x}\left( x^{n} \right) = n x^{n-1}
\]

مشتق جزئی:

\[
\pderiv{f}{x} = \lim_{h \to 0} \frac{f(x+h, y) - f(x, y)}{h}
\]

مشتق دوم:

\[
\deriv{^{2}f}{x^{2}} = \deriv{}{x}\left( \deriv{f}{x} \right)
\]

\section{انتگرال}
\label{sec:calculus-integrals}

انتگرال نامعین:

\[
\int x^{n} \dd{x} = \frac{x^{n+1}}{n+1} + C
\]

انتگرال معین:

\[
\int_{0}^{\infty} e^{-x^{2}} \dd{x} = \frac{\sqrt{\pi}}{2}
\]

انتگرال دوگانه:

\[
\iint_{D} f(x, y) \dd{x}\dd{y}
\]

\section{حد و سری}
\label{sec:calculus-limits}

حد معروف:

\[
\lim_{n \to \infty} \left(1 + \frac{1}{n}\right)^{n} = e
\]

سری تیلور:

\[
e^{x} = \sum_{n=0}^{\infty} \frac{x^{n}}{n!}
= 1 + x + \frac{x^{2}}{2!} + \frac{x^{3}}{3!} + \cdots
\]

سری فوریه:

\[
f(x) = \frac{a_{0}}{2} + \sum_{n=1}^{\infty}
\left[ a_{n} \cos(nx) + b_{n} \sin(nx) \right]
\]

%========================================================================
% CHAPTER 6 — SPECIAL FUNCTIONS
%========================================================================
\chapter{توابع ویژه}
\label{chap:special-functions}

\section{تابع گاما}
\label{sec:special-gamma}

\[
\Gamma(z) = \int_{0}^{\infty} t^{z-1} e^{-t} \dd{t}, \quad \Re(z) > 0
\]

\section{تابع سینک}
\label{sec:special-sinc}

\[
\sinc(x) = \frac{\sin(\pi x)}{\pi x}
\]

\section{تابع خطا}
\label{sec:special-erf}

\[
\erf(x) = \frac{2}{\sqrt{\pi}} \int_{0}^{x} e^{-t^{2}} \dd{t}
\]

%========================================================================
% CHAPTER 7 — NUMBER SETS
%========================================================================
\chapter{مجموعه‌های عددی}
\label{chap:number-sets}

\section{نمادهای استاندارد}
\label{sec:number-sets-symbols}

جدول \cref{tab:number-sets} مجموعه‌های عددی استاندارد را
نشان می‌دهد.

\begin{matintable}{مجموعه‌های عددی استاندارد}{tab:number-sets}
\begin{tblr}{
    width = \textwidth,
    colspec = {c l X[c]},
    hlines, vlines,
    colsep = 8pt,
    rowsep = 4pt,
    row{1} = {font=\bfseries},
}
نماد & نام & توضیحات \\
$\N$ & اعداد طبیعی & $\{1, 2, 3, \ldots\}$ \\
$\Z$ & اعداد صحیح & $\{\ldots, -2, -1, 0, 1, 2, \ldots\}$ \\
$\Q$ & اعداد گویا & $\{p/q \mid p, q \in \Z, q \neq 0\}$ \\
$\R$ & اعداد حقیقی & $(-\infty, +\infty)$ \\
$\C$ & اعداد مختلط & $\{a + bi \mid a, b \in \R\}$ \\
\end{tblr}
\end{matintable}

\section{رابطه بین مجموعه‌ها}
\label{sec:number-sets-relations}

\[
\N \subset \Z \subset \Q \subset \R \subset \C
\]

%========================================================================
% CHAPTER 8 — SUMMARY
%========================================================================
\chapter{خلاصه‌ی نتایج}
\label{chap:summary}

\section{وضعیت سیستم ریاضی}
\label{sec:summary-math-status}

جدول \cref{tab:math-status} وضعیت قابلیت‌های ریاضی
\lr{MatinBook} را نشان می‌دهد.

\begin{matintable}{وضعیت قابلیت‌های ریاضی}{tab:math-status}
\begin{tblr}{
    width = \textwidth,
    colspec = {l X[c] X[c]},
    hlines, vlines,
    colsep = 8pt,
    rowsep = 4pt,
    row{1} = {font=\bfseries},
}
بخش & قابلیت‌ها & وضعیت \\
معادلات & \lr{align}, \lr{gather}, \lr{equation} & \textcolor{matingreen}{فعال} \\
ماتریس‌ها & \lr{mat}, \lr{pmat}, \lr{vmat} & \textcolor{matingreen}{فعال} \\
دلیمترها & \lr{abs}, \lr{norm}, \lr{ceil}, \lr{floor}, \lr{inner}, \lr{set}, \lr{paren} & \textcolor{matingreen}{فعال} \\
عملگرها & \lr{grad}, \lr{curl}, \lr{diver}, \lr{rank}, \lr{trace}, \lr{diag} & \textcolor{matingreen}{فعال} \\
مشتق و انتگرال & \lr{deriv}, \lr{pderiv}, \lr{dd}, \lr{D} & \textcolor{matingreen}{فعال} \\
مجموعه‌ها & \lr{N}, \lr{Z}, \lr{Q}, \lr{R}, \lr{C} & \textcolor{matingreen}{فعال} \\
توابع ویژه & \lr{sinc}, \lr{sgn}, \lr{gcdop}, \lr{lcm}, \lr{erf} & \textcolor{matingreen}{فعال} \\
\end{tblr}
\end{matintable}

\section{معیارهای ارزیابی}
\label{sec:summary-criteria}

در این آزمون، سیستم ریاضی \lr{MatinBook} بر اساس چهار
معیار ارزیابی شد:

\begin{enumerate}
    \item \textbf{دقت نمادگذاری}: نمایش درست نمادهای پیچیده.
    \item \textbf{فاصله‌گذاری}: فاصله‌ی مناسب حول عملگرها.
    \item \textbf{شماره‌گذاری}: شماره‌گذاری منظم و پیوسته.
    \item \textbf{ارجاع هوشمند}: ارجاع دقیق با \lr{\textbackslash cref}.
\end{enumerate}

\section{نتیجه‌گیری}
\label{sec:summary-conclusion}

\textbf{تمام زیرسیستم‌های ریاضی \lr{MatinBook} نسخه‌ی
\lr{\matinbookversion} با موفقیت بارگذاری و آزمایش شدند.
معادلات، ماتریس‌ها، دلیمترها، عملگرها، مشتق و انتگرال،
توابع ویژه و مجموعه‌های عددی همگی به‌درستی کار می‌کنند.}

%------------------------------------------------------------------------
% BACK MATTER (symmetric margins)
%------------------------------------------------------------------------
\backmatter
\frontmattergeometry

% Bibliography (empty placeholder)
\printbibliography[title={منابع و مراجع}]

% Index
\printindex

%------------------------------------------------------------------------
% BACK COVER
%------------------------------------------------------------------------
\authorbio{%
    تیم توسعه‌ی \lr{MatinBook}، گروهی از پژوهشگران و برنامه‌نویسان
    فارسی‌زبان است که با هدف ارائه‌ی یک راه‌حل حرفه‌ای برای
    حروف‌چینی کتاب‌های فنی فارسی فعالیت می‌کند. این پروژه حاصل
    سال‌ها تجربه در حوزه‌ی نشر دیجیتال و برنامه‌نویسی است.
}

\makebackcover{%
    این سند، پنجمین آزمون از مجموعه‌ی آزمون‌های یکپارچگی
    \lr{MatinBook} نسخه‌ی \lr{\matinbookversion} است.

    \vspace{0.5cm}

    \textbf{زیرسیستم‌های آزموده‌شده:}
    \begin{itemize}
        \item معادلات درون‌خطی و نمایشی
        \item محیط‌های چندخطی
        \item ماتریس‌ها و بردارها
        \item دلیمترهای هوشمند
        \item عملگرهای سفارشی
        \item مشتق و انتگرال
        \item توابع ویژه
        \item مجموعه‌های عددی
    \end{itemize}
}

\end{document}